\section{Optimizer scaling：StepFun、Muon 与 scale dependence}


\singleslide{slides-images/slide-031.jpg}{Optimizer scaling 路线：optimizer choice 与 tuning 都可能随规模变化。}{31}
比较 optimizer 不能只复用同一 LR/batch；每个方法需要公平的规模相关调参和等 compute 对照。
\begin{knowledgebox}{讲义补充：源 Slide 31 的问题}
Optimizer 不是“选 AdamW 就结束”。不同 optimizer 的最优 LR、batch、weight decay、momentum 可能随 scale 变化；如果只在小模型上比较，可能因为没调好超参而误判算法。
\end{knowledgebox}

\begin{figure}[H]
\centering
\includegraphics[width=0.86\textwidth]{slide-032.jpg}
\caption{Slide 32：StepFun scaling law study，核心问题是随着 scale 如何设置 LR/batch。}
\figsource{CS336 Spring 2026 Lecture 11 官方幻灯片。}
\end{figure}

\begin{knowledgebox}{读图：Slide 32 的 StepFun 角色}
StepFun 类研究把 optimizer hyperparameter scaling 当成主要对象，而不只是模型大小。它类似 DeepSeek/Qwen 路线：训练大量小模型，在 LR/batch 空间中找规律。
\end{knowledgebox}


\singleslide{slides-images/slide-033.jpg}{LR 与 batch scaling 的变量选择：critical-batch(loss)、compute power law，或其他函数形式。}{33}
自变量决定外推语义；用 loss 表示训练状态通常比只用 step 更可迁移，但仍要验证不同模型族。
\begin{knowledgebox}{讲义补充：源 Slide 33 的变量选择}
Critical batch 把 batch 写成 loss 的函数；DeepSeek 风格可能把超参写成 compute 的幂函数；也可能与数据量、模型大小分别相关。函数形式不是数学细节，而是决定外推是否可信的建模选择。
\end{knowledgebox}

\begin{figure}[H]
\centering
\includegraphics[width=0.86\textwidth]{slide-034.jpg}
\caption{Slide 34：StepFun approach 是纯经验 grid search，训练模型来描绘 hyperparameter space。}
\figsource{CS336 Spring 2026 Lecture 11 官方幻灯片。}
\end{figure}

\begin{knowledgebox}{读图：Slide 34 的 grid search 代价}
纯经验方法的优势是不依赖强理论假设；劣势是成本高，且容易受数据范围影响。它要求每个模型规模和数据量上都有足够覆盖，否则拟合出的 scaling rule 只是局部规律。
\end{knowledgebox}

\begin{figure}[H]
\centering
\includegraphics[width=0.86\textwidth]{slide-035.jpg}
\caption{Slide 35：Observation 1：pre-training loss over batch/LR 通常呈凸形，minimizer 可较清楚识别。}
\figsource{CS336 Spring 2026 Lecture 11 官方幻灯片。}
\end{figure}

\begin{knowledgebox}{读图：Slide 35 为什么重要}
若 loss surface 对 LR/batch 近似凸，hyperparameter tuning 就可被系统化：每个 scale 找最小点，再拟合最小点随 scale 的轨迹。若曲面多峰或噪声大，外推就危险得多。
\end{knowledgebox}

\begin{figure}[H]
\centering
\includegraphics[width=0.86\textwidth]{slide-036.jpg}
\caption{Slide 36：Observation 2：batch 主要依赖 dataset size；固定模型时更大数据可能需要更高 LR，但这对 WSD 可能脆弱。}
\figsource{CS336 Spring 2026 Lecture 11 官方幻灯片。}
\end{figure}

\begin{warningbox}{读图：Slide 36 的 fragile LR 结论}
如果换成 WSD schedule、不同 warmup、不同 decay 或不同 optimizer，LR scaling 结论可能改变。Batch scaling 可能更稳定，LR scaling 更依赖训练 recipe。因此公开 LR 幂律要带着 recipe 一起读。
\end{warningbox}

\begin{figure}[H]
\centering
\includegraphics[width=0.86\textwidth]{slide-037.jpg}
\caption{Slide 37：Observation 3：结果可能泛化到 MoE 和其他 datasets，但仍需谨慎。}
\figsource{CS336 Spring 2026 Lecture 11 官方幻灯片。}
\end{figure}

\begin{knowledgebox}{读图：Slide 37 的 robustness 问题}
Scaling rule 是否跨模型族、数据集、dense/MoE 架构泛化，是 optimizer scaling 的核心。若泛化好，小模型 tuning 价值很大；若不泛化，每个新数据集/架构都要重新 sweep。
\end{knowledgebox}


\begin{importantbox}{本节核心观点}
评估 optimizer 时，必须同时控制 compute、Chinchilla ratio、batch/LR scaling 和模型规模。否则一个 optimizer 看似更好，可能只是因为它在某个 scale 上超参调得更合适。
\end{importantbox}

\begin{figure}[H]
\centering
\includegraphics[width=0.86\textwidth]{slide-038.jpg}
\caption{Slide 38：Optimizer and scale，optimization 是 LLM 的核心，但因 scale dependence 很棘手。}
\figsource{CS336 Spring 2026 Lecture 11 官方幻灯片。}
\end{figure}

\singleslide{slides-images/slide-039.jpg}{Optimizer comparison 陷阱：不同 optimizer 有不同最优超参和 scaling，默认配置会制造假赢家。}{39}
公平比较应为每个规模分别调 LR、batch、momentum/weight decay，并报告 tune budget。
\begin{warningbox}{讲义补充：源 Slide 39 的 optimizer 比较陷阱}
若 AdamW 调得很好、Muon 或 SGD 没调好，比较没有意义。每个 optimizer 都需要自己的 learning-rate/batch/weight-decay scaling rule。公平比较不是同一套超参，而是各自最优超参。
\end{warningbox}

\begin{figure}[H]
\centering
\includegraphics[width=0.86\textwidth]{slide-040.jpg}
\caption{Slide 40：Problem 2：显著 scale dependence；算法开发必须检查 compute 和 Chinchilla ratios。}
\figsource{CS336 Spring 2026 Lecture 11 官方幻灯片。}
\end{figure}

\begin{knowledgebox}{读图：Slide 40 对算法论文的提醒}
一个 optimizer 在小模型或非 Chinchilla ratio 下胜出，不代表在大模型 compute-optimal setting 中胜出。算法开发需要 scaling with compute，而不是只报告一个固定规模实验。
\end{knowledgebox}

\singleslide{slides-images/slide-041.jpg}{看似平滑的 hyperparameter scaling 仍可能在更大规模 blow up。}{41}
外推曲线需要在未参与拟合的中间规模做验证，并监控稳定性而非只看最终 loss；parametrization 或 optimizer 细节可让漂亮规律失效。
\begin{warningbox}{讲义补充：源 Slide 41 的 blow-up 教训}
Scaling curves 在小范围内平滑，不代表在更大 compute 上稳定。某些参数化、batch-size LR scaling 或 optimizer tweak 可能在大 scale 才暴露不稳定。外推前必须留出中间 scale 验证点。
\end{warningbox}

\subsection{Muon}

前面的 StepFun 结果说明 optimizer 比较必须随 scale 调参，本节用 Muon 检查一个更激进的矩阵更新方法。核心问题不是 Newton--Schulz 形式本身是否漂亮，而是它在相同 compute、数据配比和各自最优超参下，能否稳定跨越 NanoGPT、小规模 study 与大型 MoE 训练。

\begin{figure}[H]
\centering
\includegraphics[width=0.86\textwidth]{slide-042.jpg}
\caption{Slide 42：Muon 是面向 matrix-valued parameters 的 optimizer，用 Newton-Schulz 近似正交化更新矩阵。}
\figsource{CS336 Spring 2026 Lecture 11 官方幻灯片。}
\end{figure}

\begin{knowledgebox}{读公式：Slide 42 的 Muon 直觉}
若更新矩阵 \(B_t=USV^\top\)，Muon 近似把它转成 \(UV^\top\)，也就是去掉奇异值尺度，只保留方向上的正交化结构。它试图让矩阵参数的更新更均衡，但实际收益必须看 scaling 和超参调优。
\end{knowledgebox}

\begin{figure}[H]
\centering
\includegraphics[width=0.86\textwidth]{slide-043.jpg}
\caption{Slide 43：Muon and scaling：NanoGPT speedrun、小规模 scaling study、Kimi K2 都显示 Muon 可在 scale 上工作。}
\figsource{CS336 Spring 2026 Lecture 11 官方幻灯片。}
\end{figure}

\begin{knowledgebox}{读图：Slide 43 的保守结论}
Muon “works at scale” 不等于它在所有 scale、所有模型、所有 compute ratios 下都优于 AdamW。Scaling gains 很难测，因为要公平调参、控制 compute、控制数据和架构。课程更看重方法论：optimizer 改动必须做 scaling 验证。
\end{knowledgebox}

\subsection{本章小结}

Optimizer scaling 的难点在于：不同 optimizer 有不同最优超参，而且这些超参随 scale、数据、compute 和训练 schedule 变化。任何 optimizer 结论都必须绑定 scaling study，而不是孤立单点。

